% Options for packages loaded elsewhere
\PassOptionsToPackage{unicode}{hyperref}
\PassOptionsToPackage{hyphens}{url}
%
\documentclass[
  ignorenonframetext,
]{beamer}
\usepackage{pgfpages}
\setbeamertemplate{caption}[numbered]
\setbeamertemplate{caption label separator}{: }
\setbeamercolor{caption name}{fg=normal text.fg}
\beamertemplatenavigationsymbolsempty
% Prevent slide breaks in the middle of a paragraph
\widowpenalties 1 10000
\raggedbottom
\setbeamertemplate{part page}{
  \centering
  \begin{beamercolorbox}[sep=16pt,center]{part title}
    \usebeamerfont{part title}\insertpart\par
  \end{beamercolorbox}
}
\setbeamertemplate{section page}{
  \centering
  \begin{beamercolorbox}[sep=12pt,center]{part title}
    \usebeamerfont{section title}\insertsection\par
  \end{beamercolorbox}
}
\setbeamertemplate{subsection page}{
  \centering
  \begin{beamercolorbox}[sep=8pt,center]{part title}
    \usebeamerfont{subsection title}\insertsubsection\par
  \end{beamercolorbox}
}
\AtBeginPart{
  \frame{\partpage}
}
\AtBeginSection{
  \ifbibliography
  \else
    \frame{\sectionpage}
  \fi
}
\AtBeginSubsection{
  \frame{\subsectionpage}
}
\usepackage{amsmath,amssymb}
\usepackage{iftex}
\ifPDFTeX
  \usepackage[T1]{fontenc}
  \usepackage[utf8]{inputenc}
  \usepackage{textcomp} % provide euro and other symbols
\else % if luatex or xetex
  \usepackage{unicode-math} % this also loads fontspec
  \defaultfontfeatures{Scale=MatchLowercase}
  \defaultfontfeatures[\rmfamily]{Ligatures=TeX,Scale=1}
\fi
\usepackage{lmodern}
\ifPDFTeX\else
  % xetex/luatex font selection
\fi
% Use upquote if available, for straight quotes in verbatim environments
\IfFileExists{upquote.sty}{\usepackage{upquote}}{}
\IfFileExists{microtype.sty}{% use microtype if available
  \usepackage[]{microtype}
  \UseMicrotypeSet[protrusion]{basicmath} % disable protrusion for tt fonts
}{}
\makeatletter
\@ifundefined{KOMAClassName}{% if non-KOMA class
  \IfFileExists{parskip.sty}{%
    \usepackage{parskip}
  }{% else
    \setlength{\parindent}{0pt}
    \setlength{\parskip}{6pt plus 2pt minus 1pt}}
}{% if KOMA class
  \KOMAoptions{parskip=half}}
\makeatother
\usepackage{xcolor}
\newif\ifbibliography
\setlength{\emergencystretch}{3em} % prevent overfull lines
\providecommand{\tightlist}{%
  \setlength{\itemsep}{0pt}\setlength{\parskip}{0pt}}
\setcounter{secnumdepth}{-\maxdimen} % remove section numbering
\usepackage{arev} % order matters, putting this above allows FiraSans to override it for body text
\usepackage[sfdefault]{FiraSans}
\usepackage{inconsolata}
\usepackage{graphicx}
\usepackage{hyperref}
\usepackage{listings}
\usepackage{tikz}
\usepackage{pgfplots}
\usepackage{amsmath}
\usepackage[most]{tcolorbox}

\usepackage{xcolor}
\definecolor{blackish}{RGB}{56,58,54}
\definecolor{redish}{RGB}{109,41,49}
\definecolor{red}{RGB}{152,41,50}
\definecolor{orangeish}{RGB}{188,71,0}
\definecolor{blueish}{RGB}{25,33,139}

\usetheme{default}
\setbeamertemplate{navigation symbols}{}
\setbeamertemplate{itemize items}[triangle]
\setbeamertemplate{enumerate items}[default]

\setbeamercolor{structure}{fg=blueish}
\setbeamercolor{title}{fg=redish}
\setbeamercolor{subtitle}{fg=blackish}
\setbeamercolor{frametitle}{fg=blackish}
\setbeamercolor{frametitle rule}{fg=redish}
\setbeamercolor{footline}{fg=blackish}
\setbeamercolor{author}{fg=blackish}
\setbeamercolor{date}{fg=blackish}
\setbeamercolor{normal text}{fg=blackish}
\setbeamercolor{block title}{fg=white,bg=blueish!80}
\setbeamercolor{block body}{bg=blueish!5}
\setbeamercolor{block title alerted}{fg=white,bg=red!80}
\setbeamercolor{block body alerted}{bg=red!5}
\setbeamercolor{block title example}{fg=white,bg=orangeish!80}
\setbeamercolor{block body example}{bg=orangeish!5}

\setbeamertemplate{frametitle}{
  \vskip0.2em
  {\usebeamerfont{frametitle}\usebeamercolor[fg]{frametitle}\insertframetitle\par}
  {\color{redish}\hrule height 0.5pt \vspace{0.1em}}
  \vskip0.2em
}

\setbeamertemplate{footline}{%
  \hfill%
  {\color{red}\scriptsize COMP3000 Programming Languages}\hspace{0.4em}%
  {
\includegraphics[height=0.4cm]{logo.jpg}}\hspace{0.4em}%
  {\color{blackish}\scriptsize\insertframenumber/\inserttotalframenumber}\hspace{0.4em}\vspace{0.2em}
}

\setbeamertemplate{title page}{%
  \vbox{}
  \vfill
  \begingroup
    \centering
    {\usebeamerfont{title}\usebeamercolor[fg]{title}\inserttitle\par}
    \vskip0.5em
    {\usebeamerfont{subtitle}\usebeamercolor[fg]{subtitle}\insertsubtitle\par}
    \vskip1em
    {\color{redish}\hrule height 0.8pt\par}
    \vskip1em
    {\usebeamerfont{author}\usebeamercolor[fg]{author}\insertauthor\par}
    \vskip0.3em
    {\usebeamerfont{date}\usebeamercolor[fg]{date}\insertdate\par}
  \endgroup
  \vfill
}

\DeclareOldFontCommand{\bf}{\normalfont\bfseries}{\mathbf}
\DeclareOldFontCommand{\tt}{\normalfont\ttfamily}{\texttt}

\lstset{basicstyle=\ttfamily}

\newtcolorbox{aside}[1][]{
  title=Aside,
  width=0.9\textwidth,
  fonttitle=\bfseries,
  breakable,
  fonttitle=\bfseries\color{black},
  colframe=blueish!80,
  colback=blueish!2
  #1}

\newtcolorbox{goodtoknow}[1][]{
  title=Good to know,
  width=0.9\textwidth,
  fonttitle=\bfseries,
  breakable,
  fonttitle=\bfseries\color{black},
  colframe=blueish!80,
  colback=blueish!2
  #1}

\newtcolorbox{notebox}[1][]{
  title=Note,
  width=\textwidth,
  fonttitle=\bfseries,
  breakable,
  fonttitle=\bfseries\color{black},
  colframe=orangeish!80,
  colback=orangeish!2
  #1}

\newtcolorbox{hint}[1][]{
  title=Hint,
  width=\textwidth,
  fonttitle=\bfseries,
  breakable,
  fonttitle=\bfseries\color{white},
  colframe=blueish!80,
  colback=blueish!2
  #1}

\newtcolorbox{todo}[1][]{
  title=!! TODO !!,
  width=\textwidth,
  fonttitle=\bfseries,
  breakable,
  fonttitle=\bfseries\color{white},
  colframe=red!80,
  colback=red!2
  #1}

\providecommand{\tightlist}{%
  \setlength{\itemsep}{0pt}\setlength{\parskip}{0pt}}

\hypersetup{
  colorlinks=true,
  linkcolor=blueish,
  urlcolor=blueish,
  filecolor=blueish
}

\ifLuaTeX
  \usepackage{selnolig}  % disable illegal ligatures
\fi
\usepackage{bookmark}
\IfFileExists{xurl.sty}{\usepackage{xurl}}{} % add URL line breaks if available
\urlstyle{same}
\hypersetup{
  pdftitle={COMP3000 Week 2},
  pdfauthor={COMP3000 Programming Languages},
  hidelinks,
  pdfcreator={LaTeX via pandoc}}

\title{COMP3000 Week 2}
\subtitle{Little Languages + The Pipeline}
\author{COMP3000 Programming Languages}
\date{}

\begin{document}
\frame{\titlepage}

\begin{frame}{Last week in 60 seconds}
\phantomsection\label{last-week-in-60-seconds}
\begin{itemize}
\tightlist
\item
  A \textbf{domain-specific language} (DSL) is a little language built
  for one job
\item
  Examples: SQL, regex, LaTeX, Make, CSS, shell
\item
  When you can, \textbf{reuse an existing one} instead of rolling your
  own
\item
  The big idea: you will \textbf{build your own} little language for
  water flows
\item
  The pipeline, teased: \textbf{Scanner \(\rightarrow\) Parser
  \(\rightarrow\) Interpreter}
\end{itemize}

\note{Recap from Week 1. Today we turn the three-word pipeline teaser
into the full picture, and we use regex as a worked example of a little
language.}
\end{frame}

\begin{frame}{From teaser to full picture}
\phantomsection\label{from-teaser-to-full-picture}
\begin{itemize}
\tightlist
\item
  Last week: just the names -- Scanner, Parser, Interpreter
\item
  The real pipeline has more stages than that
\item
  Not every implementation walks the whole route
\item
  Today: the \textbf{complete map}, then where our implementation stops
\end{itemize}

\note{Set up the arc of the lecture: the full mountain, then why
``compiler or interpreter?'' is the wrong question, then regex as a DSL
that shows every idea in miniature.}
\end{frame}

\begin{frame}{Little languages, again}
\phantomsection\label{little-languages-again}
\begin{itemize}
\tightlist
\item
  Ch 1: ``For every successful general-purpose language, there are a
  thousand successful niche ones''
\item
  They are \textbf{pidgins tailor-built to a specific task}
\item
  Small on purpose: restricted syntax, restricted features
\item
  They lean on the machinery around them (their ``host'') for the hard
  parts
\end{itemize}

\note{The definition again, but deeper. A DSL deliberately gives up
generality to speak precisely about one domain. Think of sed, or the
command language in a game: it does one job well and leans on the
surrounding environment.}
\end{frame}

\begin{frame}{Reuse\ldots{} or roll your own?}
\phantomsection\label{reuse-or-roll-your-own}
\begin{itemize}
\tightlist
\item
  ``When you can, good to reuse an existing one instead of rolling your
  own''
\item
  Why: documentation, debuggers, editor support, syntax highlighting
\item
  ``Once you factor in all the other trappings, doing it yourself
  becomes a tall order''
\item
  But when nothing fits your domain, building one is the right call --
  that's this unit
\end{itemize}

\note{The trade-off from ch 1. Our domain (water flows) has no existing
language, so we build one. But the default should always be to reuse.}
\end{frame}

\begin{frame}{The pipeline, in full}
\phantomsection\label{the-pipeline-in-full}
\begin{center}
\includegraphics[width=0.85\textwidth]{text/site/image/a-map-of-the-territory/mountain.png}
\end{center}

\textbf{Scan \(\rightarrow\) Parse \(\rightarrow\) Static Analysis
\(\rightarrow\) Optimise \(\rightarrow\) Code Gen \(\rightarrow\) VM
\(\rightarrow\) Runtime}

\note{Ch 2's map of the territory. Source text at the bottom, climb to
where meaning is clear, descend to something the machine can run. We
covered the first two last week; today we name the rest.}
\end{frame}

\begin{frame}{Front end, middle end, back end}
\phantomsection\label{front-end-middle-end-back-end}
\begin{itemize}
\tightlist
\item
  \textbf{Front end}: where source text becomes a structured
  representation

  \begin{itemize}
  \tightlist
  \item
    scanning and parsing
  \end{itemize}
\item
  \textbf{Middle end}: where we reason about the program

  \begin{itemize}
  \tightlist
  \item
    static analysis and optimisation
  \end{itemize}
\item
  \textbf{Back end}: where we produce something runnable

  \begin{itemize}
  \tightlist
  \item
    code generation, virtual machine, runtime
  \end{itemize}
\end{itemize}

\note{The three traditional regions of a compiler. Our tree-walk
interpreter only uses the front end and stops there -- it evaluates the
tree directly.}
\end{frame}

\begin{frame}{Static analysis: seeing what the code does}
\phantomsection\label{static-analysis-seeing-what-the-code-does}
\begin{itemize}
\tightlist
\item
  The parser knows the \emph{syntax}; static analysis works out the
  \emph{meaning}
\item
  \textbf{Binding}: connecting each identifier to the variable or
  function it refers to
\item
  \textbf{Type checking}: rejecting programs that don't make sense (in
  typed languages)
\item
  Ch 2: ``get a high-level view of what the code is doing''
\item
  In Lox: a \textbf{resolver} pass walks the tree before we evaluate it
\end{itemize}

\note{Static analysis runs without executing anything. It's where a lot
of errors get caught early. Our Lox interpreter adds this pass in the
book's second part.}
\end{frame}

\begin{frame}[fragile]{Optimisation: doing less work}
\phantomsection\label{optimisation-doing-less-work}
\begin{itemize}
\tightlist
\item
  If an expression always gives the same value, compute it once at build
  time
\item
  \textbf{Constant folding}:
\end{itemize}

\vspace{0.5em}

\texttt{pennyArea\ =\ 3.14159\ *\ (0.75\ /\ 2)\ *\ (0.75\ /\ 2);}

becomes

\texttt{pennyArea\ =\ 0.4417860938;}

\begin{itemize}
\tightlist
\item
  ``Optimization is a huge part of the programming language business''
\end{itemize}

\note{Constant folding from ch 2. We can do the arithmetic in the
compiler and replace the expression with its result. Later the book
pairs it with constant propagation.}
\end{frame}

\begin{frame}{Code generation: speaking the machine's language}
\phantomsection\label{code-generation-speaking-the-machines-language}
\begin{itemize}
\tightlist
\item
  The back end converts our representation into runnable code
\item
  \textbf{Native code}: instructions for a real chip (x86, ARM)

  \begin{itemize}
  \tightlist
  \item
    lightning fast, but tied to one architecture, and a lot of work
  \end{itemize}
\item
  \textbf{Bytecode}: instructions for a \emph{virtual} machine

  \begin{itemize}
  \tightlist
  \item
    ``a dense, binary encoding of the language's low-level operations''
  \item
    portable across architectures
  \end{itemize}
\end{itemize}

\note{The design decision at the bottom of the mountain. Wirth and
Richards produced p-code for a hypothetical machine so one compiler
could target many chips.}
\end{frame}

\begin{frame}{Virtual machines and the runtime}
\phantomsection\label{virtual-machines-and-the-runtime}
\begin{itemize}
\tightlist
\item
  A \textbf{VM} emulates a hypothetical chip; it runs the bytecode at
  runtime
\item
  Slower than native code, but simple and portable

  \begin{itemize}
  \tightlist
  \item
    ``Implement your VM in C, and you can run your language on any
    platform with a C compiler''
  \end{itemize}
\item
  The \textbf{runtime} provides services while a program runs:

  \begin{itemize}
  \tightlist
  \item
    garbage collection, and tracking object types at runtime
  \end{itemize}
\end{itemize}

\note{Language VMs (like Java's JVM, or Python's) vs system VMs (like
VirtualBox). Our second interpreter in this unit is a bytecode VM.}
\end{frame}

\begin{frame}{Where our journey stops}
\phantomsection\label{where-our-journey-stops}
\begin{itemize}
\tightlist
\item
  Our first interpreter is a \textbf{tree-walk interpreter}: scan,
  parse, evaluate

  \begin{itemize}
  \tightlist
  \item
    no static analysis, no code generation, no VM -- it walks the tree
    directly
  \end{itemize}
\item
  Later in the unit we add the static-analysis pass
\item
  So we climb part of the mountain, and evaluate from the top
\end{itemize}

\note{Reassure students: we don't build the whole mountain. We take the
tree-walk route, and the book's second half adds a bytecode VM. The
other stages are context.}
\end{frame}

\begin{frame}{Compilers and interpreters}
\phantomsection\label{compilers-and-interpreters}
\begin{itemize}
\tightlist
\item
  ``What's the difference?'' is like asking the difference between a
  fruit and a vegetable
\item
  \textbf{Fruit} is a \emph{botanical} term; \textbf{vegetable} is
  \emph{culinary}
\item
  Neither implies the negation of the other
\item
  Tomatoes are both; apples are fruit but not vegetables; carrots are
  vegetables but not fruit
\end{itemize}

\begin{center}
\includegraphics[width=0.55\textwidth]{text/site/image/a-map-of-the-territory/plants.png}
\end{center}

\note{The ch 2 analogy. ``Compiled'' and ``interpreted'' are not two
sides of a binary -- they are different ways of describing an
implementation, and both can apply at once.}
\end{frame}

\begin{frame}{Compiling is an implementation technique}
\phantomsection\label{compiling-is-an-implementation-technique}
\begin{itemize}
\tightlist
\item
  \textbf{Compiling}: translating source to some other, usually
  lower-level, form

  \begin{itemize}
  \tightlist
  \item
    bytecode, machine code, or even another high-level language
    (transpiling)
  \end{itemize}
\item
  A \textbf{compiler} translates source but doesn't execute it -- you
  run the output yourself
\item
  An \textbf{interpreter} takes source and executes it directly, ``from
  source''
\end{itemize}

\note{Three crisp definitions from ch 2. These are properties of
\emph{implementations}, not of languages.}
\end{frame}

\begin{frame}[fragile]{It's not about the language}
\phantomsection\label{its-not-about-the-language}
\begin{itemize}
\tightlist
\item
  \textbf{CPython}: runs your program from source (an interpreter), and
  it \emph{has} a compiler -- it parses and compiles to bytecode
  internally
\item
  \textbf{Go}: \texttt{go\ build} compiles and stops; \texttt{go\ run}
  compiles \emph{then} executes -- it both \emph{is} and \emph{has} a
  compiler
\item
  So ``is Java compiled or interpreted?'' is not meaningful
\item
  We can only answer for \textbf{particular implementations}
\end{itemize}

\begin{center}
\includegraphics[width=0.6\textwidth]{text/site/image/a-map-of-the-territory/venn.png}
\end{center}

\note{The overlapping region of the Venn diagram: interpreters that
compile internally. Our second interpreter lives there too. Most
``scripting'' languages work this way.}
\end{frame}

\begin{frame}[fragile]{Regex: a DSL for text}
\phantomsection\label{regex-a-dsl-for-text}
\begin{itemize}
\tightlist
\item
  \textbf{Regular expressions} are a little language for pattern
  matching in text
\item
  Match a pattern against a string,
  e.g.~\texttt{\textbackslash{}d\{4\}-\textbackslash{}d\{2\}-\textbackslash{}d\{2\}}
  matches a date
\item
  Look at its features and it is suspiciously like a general-purpose
  language:

  \begin{itemize}
  \tightlist
  \item
    \textbf{variables}: capture groups, reused later
  \item
    \textbf{loops}: the \texttt{*} and \texttt{+} quantifiers
  \item
    \textbf{conditions}: the \texttt{?} operator\ldots{} sort of
  \end{itemize}
\end{itemize}

\note{Regex is the worked DSL of the lecture. The features list here
sets up the completeness discussion at the end.}
\end{frame}

\begin{frame}{Live on regexr.com}
\phantomsection\label{live-on-regexr.com}
\begin{itemize}
\tightlist
\item
  Head to \textbf{regexr.com} and try it yourself
\item
  Try:

  \begin{itemize}
  \tightlist
  \item
    an exact string, e.g.~``my presence''
  \item
    one of two strings, e.g.~``my presence'' or ``my friends''
  \item
    a repeated letter, e.g.~repeated p's or g's
  \item
    capture the first word of each character's dialogue
  \end{itemize}
\end{itemize}

\note{Live demo on regexr.com. This mirrors the Week 2 workshop
activity, where you do the same exploration and then write up your own
example in the FAT.}
\end{frame}

\begin{frame}{What makes a language complete?}
\phantomsection\label{what-makes-a-language-complete}
\begin{itemize}
\tightlist
\item
  A language is \textbf{complete} if it can express all possible
  computations
\item
  All complete languages have equivalent expressive power
\item
  The \textbf{Turing view}: variables, conditions, and loops -- like a
  Turing machine
\item
  The \textbf{Church view}: variables and functions -- like the lambda
  calculus
\item
  ``This dual nature of computing is one of the most fascinating aspects
  of programming language theory''
\end{itemize}

\note{The completeness test from the workshop. Knowing for sure takes
real maths, but these two feature tests are a good approximation.}
\end{frame}

\begin{frame}{Wrap-up and Week 3 prep}
\phantomsection\label{wrap-up-and-week-3-prep}
\begin{itemize}
\tightlist
\item
  Today's reading: \textbf{Chapters 1--2} (re-read if needed)
\item
  \textbf{Before next class}: read \textbf{Chapter 3} and complete the
  Week 3 iRAT
\item
  Next week: \textbf{Lox} -- the language we build, and its features
\item
  Preview: a Lox program that drives turtle graphics
\end{itemize}

\note{Next week introduces Lox and the first real implementation work.
Chapter 3 is short -- the iRAT will draw directly on it.}
\end{frame}

\end{document}
